\documentclass[a4paper, 12pt]{article}
\usepackage[utf8]{inputenc} % Encodage UTF-8
\usepackage[T1]{fontenc}    % Encodage des caractères
\usepackage[francais]{babel} % Langue française
\usepackage{geometry}       % Mise en page
\usepackage{graphicx}       % Inclusion d'images
\usepackage{hyperref}       % Hyperliens
\usepackage{fancyhdr}       % En-têtes et pieds de page
\usepackage{titlesec}       % Personnalisation des sections
\usepackage{enumitem}       % Personnalisation des listes

% Configuration de la mise en page
\geometry{margin=2.5cm}
\pagestyle{fancy}
\fancyhf{}
\fancyhead[L]{\textbf{Compte Rendu}} % Titre en en-tête gauche
\fancyhead[R]{\thepage}             % Numéro de page à droite
\fancyfoot[L]{Nom du projet}        % Pied de page gauche
\fancyfoot[R]{\today}               % Date à droite

% Configuration des sections
\titleformat{\section}{\large\bfseries}{\thesection.}{1em}{}
\titleformat{\subsection}{\normalsize\bfseries}{\thesubsection.}{1em}{}
\titleformat{\subsubsection}{\small\bfseries}{\thesubsubsection.}{1em}{}

\setlist[itemize]{label=--, left=0cm}

\begin{document}

% Page de titre
\begin{titlepage}
    \centering
    \includegraphics[width=0.2\textwidth]{logo.png} \\ % Logo (à personnaliser)
    \vspace{1cm}
    {\Large \textbf{Compte Rendu}} \\[0.5cm]
    {\large \textbf{Nom du projet ou sujet}} \\[1.5cm]
    \begin{tabbing}
        \hspace{3cm} \= \textbf{Date :} \hspace{0.5cm} \= \today \\[0.2cm]
        \> \textbf{Responsable :} \> Nom du responsable \\[0.2cm]
        \> \textbf{Participants :} \> Liste des participants \\[0.2cm]
        \> \textbf{Lieu :} \> Lieu de la réunion \\[0.2cm]
        \> \textbf{Durée :} \> Durée approximative
    \end{tabbing}
    \vfill
\end{titlepage}

% Introduction
\section*{Introduction}
Décrivez brièvement l'objectif de la réunion ou du projet, ainsi que les points importants à traiter.

% Ordre du jour
\section*{Ordre du jour}
\begin{itemize}
    \item Point 1 : Description du point.
    \item Point 2 : Description du point.
    \item Point 3 : Description du point.
\end{itemize}

% Développement
\section{Déroulement et discussions}
\subsection{Point 1 : [Nom du point]}
Résumé des discussions, décisions prises, et points à suivre.

\subsection{Point 2 : [Nom du point]}
Résumé des discussions, décisions prises, et points à suivre.

\subsection{Point 3 : [Nom du point]}
Résumé des discussions, décisions prises, et points à suivre.

% Actions à entreprendre
\section{Actions à entreprendre}
\begin{itemize}
    \item \textbf{Action 1 :} Description, responsable, et échéance.
    \item \textbf{Action 2 :} Description, responsable, et échéance.
    \item \textbf{Action 3 :} Description, responsable, et échéance.
\end{itemize}

% Conclusion
\section*{Conclusion}
Résumé global, rappel des actions prioritaires, et prochaines étapes (si applicable).

% Annexe (optionnel)
\section*{Annexes}
Ajoutez des documents supplémentaires, graphiques ou détails techniques si nécessaire.

\end{document}